\bigskip
\noindent Autres pistes de l'onglet Findings pour ce thème, dans des cotes qui n'ont pas de section ici :\par

\begin{encadre}
\begin{conjecture*}[Grothendieck, pages 37 et 41, algèbres bivariantes simpliciales et algèbres de Poincaré\piste{11-simplicial-bivariant-poincare-equivalence}]
L'évaluation en un ensemble à un élément est une équivalence entre les algèbres bivariantes simpliciales, foncteurs contravariants des ensembles finis non vides vers les algèbres graduées bivariantes de la page~33, dont la valeur en un point est une algèbre de Poincaré, et les algèbres de Poincaré. Le foncteur $A(\cdot) \mapsto (A(\mathrm{pt}), \Delta)$, avec $\Delta = \delta_*(1)$ de degré $d(A)$, est pleinement fidèle vers les algèbres jaugées munies d'un tel~$\Delta$.
\end{conjecture*}
\end{encadre}
{\small La page laisse \q{à préciser} les axiomes sur $\Delta$. L'édition restreint l'équivalence aux objets \q{spéciaux}, sans quoi rien ne lie $A(I)$ à $A(\mathrm{pt})$ pour $\mathrm{card}\,I \geq 2$. Mathlib n'a pas ces théories bivariantes. On se reportera aux notes de Grothendieck (\lot{11}{2}{37}).\par}

\begin{encadre}
\begin{conjecture*}[Grothendieck, pages 9 à 16, théorème de Griffiths par voie $\ell$-adique\piste{36-griffiths-conditional-l-adic}]
En toute caractéristique, soient $X$ projective lisse de dimension $2n$ et un pinceau de Lefschetz, de section hyperplane générique $Y$. Supposons a) un scindage de type Lefschetz de l'inclusion de $Y$ par une correspondance algébrique, et b) qu'aucune pièce de $\mathrm{Coker}(H^{2n-1}(X) \to H^{2n-1}(Y))$ ne soit de coniveau $n-1$. Alors un cycle de $X$ dont la restriction à $Y$ est algébriquement équivalente à zéro est homologiquement nul à torsion près, et $\mathrm{Grif}(Y) \otimes \mathbf Q$ a un sous-quotient qui s'envoie sur les classes algébriques primitives de~$X$.
\end{conjecture*}
\end{encadre}
{\small Les feuillets en esquissent une démonstration incomplète (pages 9 à 15) : le n°\,4 renvoie à des \q{calculs explicites essentiellement triviaux}, et l'argument est de la lecture. La page dit \q{sous-espace} là où l'argument ne donne qu'un sous-quotient. L'hypothèse a) est un cas des conjectures standard, et b) n'est démontrée qu'en caractéristique nulle. Mathlib n'a pas la cohomologie $\ell$-adique ni l'application d'Abel--Jacobi. On se reportera aux notes de Grothendieck (\lot{36}{1}{9}).\par}

\begin{encadre}
\begin{theoreme*}[Grothendieck, pages 60 à 70, défaut des cycles invariants en $H^1$\piste{36-h1-invariant-cycle-defect}]\nonverifie
Soient $X$ projectif et régulier sur un trait strictement hensélien $S$, avec $\mathcal O_S \simeq f_*\mathcal O_X$, et $n$ premier à $p$. On a une suite exacte $0 \to {}_n\Gamma \to H^1(X, \mu_n) \to H^1(X_{\bar\eta}, \mu_n)^I \to \mathrm{Ker}(w_n) \to 0$, où $\Gamma$ est le groupe des diviseurs portés par la fibre spéciale modulo $X_0$ et $w : \Gamma \to \mathrm{NS}(X_0)$. À la limite, $H^1(X, \mathbf Z_\ell(1)) \to H^1(X_{\bar\eta}, \mathbf Z_\ell(1))^I$ est injectif de conoyau $R \otimes \mathbf Z_\ell$, où $R = \mathrm{Ker}\,w$ est cyclique d'ordre divisant le pgcd des multiplicités $d_i$ des composantes de $X_0$, et c'est un isomorphisme quand $X_\eta$ a un $0$-cycle de degré~$1$.
\end{theoreme*}
\end{encadre}
{\small Le tapuscrit des pages 60 à 63 l'énonce ; le manuscrit des pages 64 à 70 le démontre, la torsion de $\mathrm{Ker}\,w$ étant renvoyée à Raynaud et tenue \q{au moins modulo résolution des singularités}. Celle-ci est démontrée en caractéristique nulle (Hironaka) et en dimension au plus trois en toute caractéristique (Cossart--Piltant), et reste ouverte en caractéristique positive à partir de la dimension quatre ; ces trois énoncés sont formalisés dans le dépôt Formal Conjectures \cite{formalconjectures}. Mathlib n'a pas la cohomologie étale ni le groupe de Néron--Severi. On se reportera aux notes de Grothendieck (\lot{36}{3}{60}).\par}

\begin{encadre}
\begin{theoreme*}[Grothendieck, pages 73 à 76, orthogonalité sur une base adique régulière\piste{36-orthogonality-regular-adic-base}]\nonverifie
Soient $X = \operatorname{Spec} A$ une base noethérienne régulière complète pour la topologie $J$-adique, $Y = V(J)$ et $U = X - Y$. Soient $G$, $G'$ des schémas en groupes commutatifs plats sur $X$ dont les restrictions $A_U$, $A'_U$ à $U$ sont des schémas abéliens, $\xi$ une correspondance divisorielle sur $A_U \times A'_U$, $H \subset G_Y$ un sous-groupe plat à torsion $\ell^\nu$ finie et $T' \subset G'_Y$ un tore. Alors $\varphi_\xi(T_\ell(\tilde H)|_U, T_\ell(\tilde T')|_U) = 0$ en tout point de $U$ : l'orthogonalité de la partie torique et de la partie de torsion finie pour l'accouplement de Weil ne demande pas un anneau de valuation discrète.
\end{theoreme*}
\end{encadre}
{\small La réduction à un anneau de valuation discrète à corps résiduel fini et l'argument des poids de Frobenius (pages 74 à 76) sont lus en grande partie à travers des mots incertains ou illisibles, et l'équivariance galoisienne du lemme sur les corps finis est ajoutée par l'édition. SGA~7, IX, donne l'orthogonal exact sur un trait. Mathlib n'a pas les schémas abéliens ni l'accouplement de Weil. On se reportera aux notes de Grothendieck (\lot{36}{4}{73}).\par}
